% GENERATED FILE. Edit canonical lessons, then run npm run generate:books.
% canonical-source-hash: fnv1a64:d3b2808e13d413ab
% canonical-lessons: ZH-C201-youpiao, ZH-C201-kuaidi, ZH-C201-youjian, ZH-C201-fasong, ZH-C201-shoudao

\chapter{Post and Parcels}
\label{ch:zh-post-and-parcels}
\hlchaptermodality{201}

\begin{chapteropening}
\textbf{By the end of this chapter:} \emph{I can buy stamps, send things by courier, send an email and say I have got it.}

Complete the last lesson of chapter 201: I can buy stamps, send things by courier, send an email and say I have got it.
\end{chapteropening}

\section[yóupiào]{yóupiào — postage stamp}

\label{lesson:ZH-C201-youpiao}



\begin{quote}
Before the new one: say the Mandarin for to check in, then the Mandarin for to check out.
\end{quote}

\subsection*{You'll want to know: yóupiào}
\emph{yóupiào} — \textquotedblleft{}postage stamp\textquotedblright{}.

\begin{quote}
\emph{wǒ yào mǎi liǎng zhāng yóupiào} — I'd like two stamps, please.
\end{quote}

Its characters come later; for now it is a word for the ear.

\begin{cousinweb}
\emph{yóupiào} is built from \emph{yóu} 'post' and \zh{票} 'a slip'. The first is inherited from Middle Chinese \emph{hjuw} (Guangyun, in Baxter's transcription), glossed as a lodge at a border; no dictionary we use records an older form of \zh{票}, so its history is not settled.
\end{cousinweb}

\begin{grammarlens}[title={What you've built}]
One more word for reading notices and messages.
\end{grammarlens}

\subsection*{Your turn}
\begin{itemize}
\raggedright
  \item \emph{Say it:} \emph{yóupiào}
  \item \emph{Say it:} \emph{yóupiào}, once more
  \item \emph{Say it:} \emph{wǒ yào mǎi liǎng zhāng yóupiào} — I'd like two stamps, please.
  \item \emph{Recall:} say the Mandarin for to check in, then the Mandarin for to check out, then say \emph{yóupiào} again
\end{itemize}

\begin{tcolorbox}[breakable,colback=teal!4,colframe=teal!35!black,title={Before you move on}]
What does yóupiào mean? (\textquotedblleft{}Postage stamp\textquotedblright{}.) Say it once more.
\end{tcolorbox}
\section[kuàidì]{kuàidì — express delivery; courier}

\label{lesson:ZH-C201-kuaidi}



\begin{quote}
Before the new one: say the Mandarin for to check out, then the Mandarin for postage stamp.
\end{quote}

\subsection*{You'll want to know: kuàidì}
\emph{kuàidì} — \textquotedblleft{}express delivery; courier\textquotedblright{}.

\begin{quote}
\emph{wǒ yòng kuàidì jì shū} — I'm sending the books by courier.
\end{quote}

Its characters come later; for now it is a word for the ear.

\begin{cousinweb}
\emph{kuàidì} is built from \zh{快} 'fast' and \emph{dì} 'to pass along', both inherited from Middle Chinese, \emph{khwaejH} and \emph{dejX} (Guangyun, in Baxter's transcription). The dictionary glosses the first 'pleased' and the second 'to take turns'; how the pair came to mean a courier is not settled.
\end{cousinweb}

\begin{grammarlens}[title={What you've built}]
One more word for reading notices and messages.
\end{grammarlens}

\subsection*{Your turn}
\begin{itemize}
\raggedright
  \item \emph{Say it:} \emph{kuàidì}
  \item \emph{Say it:} \emph{kuàidì}, once more
  \item \emph{Say it:} \emph{wǒ yòng kuàidì jì shū} — I'm sending the books by courier.
  \item \emph{Recall:} say the Mandarin for to check out, then the Mandarin for postage stamp, then say \emph{kuàidì} again
\end{itemize}

\begin{tcolorbox}[breakable,colback=teal!4,colframe=teal!35!black,title={Before you move on}]
What does kuàidì mean? (\textquotedblleft{}Express delivery; courier\textquotedblright{}.) Say it once more.
\end{tcolorbox}
\section[yóujiàn]{yóujiàn — email; mail}

\label{lesson:ZH-C201-youjian}



\begin{quote}
Before the new one: say the Mandarin for postage stamp, then the Mandarin for courier service.
\end{quote}

\subsection*{You'll want to know: yóujiàn}
\emph{yóujiàn} — \textquotedblleft{}email; mail\textquotedblright{}.

\begin{quote}
\emph{lǎoshī de yóujiàn hěn zhòngyào} — The teacher's email is important.
\end{quote}

Its characters come later; for now it is a word for the ear.

\begin{cousinweb}
\emph{yóujiàn} is built from \emph{yóu} 'post' and \emph{jiàn} 'an item', both inherited from Middle Chinese, \emph{hjuw} and \emph{gjenX} (Guangyun, in Baxter's transcription), which glosses the second 'a portion'.
\end{cousinweb}

\begin{grammarlens}[title={What you've built}]
One more word for reading notices and messages.
\end{grammarlens}

\subsection*{Your turn}
\begin{itemize}
\raggedright
  \item \emph{Say it:} \emph{yóujiàn}
  \item \emph{Say it:} \emph{yóujiàn}, once more
  \item \emph{Say it:} \emph{lǎoshī de yóujiàn hěn zhòngyào} — The teacher's email is important.
  \item \emph{Recall:} say the Mandarin for postage stamp, then the Mandarin for courier service, then say \emph{yóujiàn} again
\end{itemize}

\begin{tcolorbox}[breakable,colback=teal!4,colframe=teal!35!black,title={Before you move on}]
What does yóujiàn mean? (\textquotedblleft{}Email; mail\textquotedblright{}.) Say it once more.
\end{tcolorbox}
\section[fāsòng]{fāsòng — to send electronically; to send a message or an email}

\label{lesson:ZH-C201-fasong}



\begin{quote}
Before the new one: say the Mandarin for courier service, then the Mandarin for email.
\end{quote}

\subsection*{You'll want to know: fāsòng}
\emph{fāsòng} — \textquotedblleft{}to send electronically; to send a message or an email\textquotedblright{}.

\begin{quote}
\emph{wǒ yǐjīng fāsòngle yóujiàn} — I've already sent the email.
\end{quote}

Its characters come later; for now it is a word for the ear.

\begin{cousinweb}
\emph{fāsòng} is built from \emph{fā} 'to send out' and \emph{sòng} 'to send off', both inherited from Middle Chinese, \emph{pjot} and \emph{suwngH} (Guangyun, in Baxter's transcription), which glosses the first 'to start' and 'to open out', and the second 'to send away'.
\end{cousinweb}

\begin{grammarlens}[title={What you've built}]
One more word for reading notices and messages.
\end{grammarlens}

\subsection*{Your turn}
\begin{itemize}
\raggedright
  \item \emph{Say it:} \emph{fāsòng}
  \item \emph{Say it:} \emph{fāsòng}, once more
  \item \emph{Say it:} \emph{wǒ yǐjīng fāsòngle yóujiàn} — I've already sent the email.
  \item \emph{Recall:} say the Mandarin for courier service, then the Mandarin for email, then say \emph{fāsòng} again
\end{itemize}

\begin{tcolorbox}[breakable,colback=teal!4,colframe=teal!35!black,title={Before you move on}]
What does fāsòng mean? (\textquotedblleft{}To send electronically; to send a message or an email\textquotedblright{}.) Say it once more.
\end{tcolorbox}
\section[shōudào]{shōudào — to get something sent; to receive}

\label{lesson:ZH-C201-shoudao}



\begin{quote}
Before the new one: say the Mandarin for email, then the Mandarin for to send electronically.
\end{quote}

\subsection*{You'll want to know: shōudào}
\emph{shōudào} — \textquotedblleft{}to get something sent; to receive\textquotedblright{}.

\begin{quote}
\emph{nǐ shōudào wǒ de yóujiàn le ma} — Did you get my email?
\end{quote}

Its characters come later; for now it is a word for the ear.

\begin{cousinweb}
\emph{shōudào} is built from \emph{shōu} 'to gather in' and \zh{到} 'to arrive', both inherited from Middle Chinese, \emph{syuw} and \emph{tawH} (Guangyun, in Baxter's transcription): taking in what has arrived.
\end{cousinweb}

\begin{grammarlens}[title={What you've built}]
One more word for reading notices and messages.
\end{grammarlens}

\subsection*{Your turn}
\begin{itemize}
\raggedright
  \item \emph{Say it:} \emph{shōudào}
  \item \emph{Say it:} \emph{shōudào}, once more
  \item \emph{Say it:} \emph{nǐ shōudào wǒ de yóujiàn le ma} — Did you get my email?
  \item \emph{Recall:} say the Mandarin for email, then the Mandarin for to send electronically, then say \emph{shōudào} again
\end{itemize}

\begin{tcolorbox}[breakable,colback=teal!4,colframe=teal!35!black,title={Before you move on}]
What does shōudào mean? (\textquotedblleft{}To get something sent; to receive\textquotedblright{}.) Say it once more.
\end{tcolorbox}
